\section{Benchmark Construction}

To ensure that validation-side optimization reliably predicts held-out evaluation performance, benchmark tasks should satisfy two key properties: \emph{Harness Sensitivity}, where task performance is sensitive to harness improvements while controlling for the underlying model, and \emph{Cross-Suite Alignment}, where validation and evaluation suites exhibit consistent performance to harness variation. To this end, we design a Two-Stage Harness-Guided Benchmark Construction Framework. As shown in Figure~\ref{fig:benchmark_construction},
we first leverage auxiliary tasks to conduct controlled evolution experiments and generate a diverse set of auxiliary evolved harnesses, and then use these harnesses to characterize task-level harness sensitivity and difficulty for constructing the final benchmark task suites.

\label{sec:task-construction}
\begin{figure}[htbp]
    \centering
    \includegraphics[width=1.0\linewidth]{Figures/benchmark_construction.pdf}
    \caption{Overview of Two-Stage Harness-Guided Benchmark Construction Framework.}
    \label{fig:benchmark_construction}
\end{figure}
% \subsection{Stage 1: Calibration Task Construction}
\subsection{Stage 1: Auxiliary Harness Generation}
\label{sec:calibration-tasks}

\paratitle{Auxiliary Tasks Collection.}
We collect auxiliary tasks from diverse sources that are corpus- and instance-level disjoint from the five source benchmark datasets of Evo-Bench, including MiroRL~\cite{2025mirorl}, RedSearcher~\cite{redsearcher2026}, Auto-ClawEval~\cite{li2026clawenvkitautomaticenvironmentgeneration}, and internally constructed datasets. 
We then evaluate these tasks with DeepSeek-V4-Flash~\cite{deepseekv4} through multiple independent rollouts and select tasks with lower average scores and longer interaction trajectories, which indicate greater room for harness improvements and stronger signals for harness evolution. This process results in a 320-task auxiliary set, comprising 128 search tasks, 128 office tasks, and 64 general agent tasks.

% \paratitle{Harnesses Construction via Evolution on Auxiliary Tasks.}
\paratitle{Auxiliary Harness Evolution.}
After constructing the auxiliary task set, we run independent evolution on it with multiple frontier models, collecting resulting harnesses throughout the process. We then deduplicate and select a representative subset to ensure sufficient diversity and coverage.

Specifically, we conduct four independent evolution experiments using four frontier models: GLM-5.2~\cite{glm52}, Claude-Opus-4.8~\cite{anthropic2026opus}, Claude-Sonnet-5~\cite{anthropic2026sonnet}, and GPT-5.6-Sol~\cite{gpt56}. Each experiment follows the same policy model, evolution protocol, and resource budget as the benchmark evaluation. We collect 73 formally evaluated harness variants and apply a deterministic diversity-aware selection procedure to identify 12 representative harnesses that maximize diversity across harness capabilities, tool orchestration, program structure, and auxiliary-task performance. We denote this selected set as $\mathcal{H}_{\mathrm{aux}}=\{h_1,\ldots,h_K\}$, where $K=12$.



\subsection{Stage 2: Harness-Guided Task Selection}
\label{sec:harness-response-prior}

We evaluate candidate tasks under the selected harness set $\mathcal{H}_{\mathrm{aux}}$ to identify tasks that reliably reflect harness improvements. Specifically, we evaluate 2,329 candidate tasks collected from APEX-Agents, BrowseComp, Claw-Eval, GDPval, and HLE using the 12 selected harnesses. We exclude multimodal tasks to maintain a unified interface and downsample BrowseComp and HLE to control evaluation cost while preserving their original task distributions.


For task $x$ from the public benchmark $\mathcal{D}_{\mathrm{public}}$, let $m_h(x)$ denote its score under harness $h\in\mathcal{H}_{\mathrm{aux}}$. 
A straightforward criterion is to select tasks with high score variance across harnesses.
However, high variance only captures the magnitude of performance differences and does not indicate whether these differences reflect the relative quality of harnesses.
% Similarly, average task scores mainly measure task difficulty and provide no information about whether a task can distinguish stronger harnesses from weaker ones.
To more precisely measure whether a task consistently responds to harness quality, we compute the correlation between the task and overall harness quality, and combine this sensitivity measure with the average task score for task selection.


\paratitle{Harness Sensitivity and Task Difficulty.}
The score of an individual task is obtained directly from its execution result, while the overall quality of each harness is measured by its average performance across the remaining tasks. We compute the Pearson correlation coefficient between task-level scores and overall harness quality as the metric for task sensitivity. 
% For each task, we further use its average score across the 12 harnesses as the difficulty measure. 
Specifically, we define two task-level metrics: harness sensitivity and average performance, as follows:

\begin{equation*}
\begin{aligned}
\mathrm{Sens}(x)
&=
\operatorname{corr}
\left(
\{m_h(x)\}_{h\in\mathcal{H}_{\mathrm{aux}}},
\{Q_h^{(-x)}\}_{h\in\mathcal{H}_{\mathrm{aux}}}
\right),\\
\mathrm{Perf}(x)
&=
\frac{1}{|\mathcal{H}_{\mathrm{aux}}|}
\sum_{h\in\mathcal{H}_{\mathrm{aux}}}m_h(x),
\end{aligned}
\end{equation*}
where $Q_h^{(-x)}$ denotes the leave-one-task-out average performance of harness $h$, computed over all tasks except $x$, and serves as a robust estimate of its overall quality.

Higher $\mathrm{Sens}(x)$ indicates that task $x$ consistently reflects the quality ranking of different harnesses, where stronger harnesses achieve better performance on the task. Therefore, tasks with higher sensitivity provide more reliable signals for evaluating harness evolution. In contrast, $\mathrm{Perf}(x)$ measures the average performance across harnesses, and we define task difficulty as $1-\mathrm{Perf}(x)$, where lower average performance indicates larger performance headroom.

% \begin{equation*}
% \begin{aligned}
%     \alpha(x)
%     &=
%     \operatorname{corr}_{h\in\mathcal{H}_{\mathrm{sel}}}
%     \bigl(m_h(x),\theta_h^{(-x)}\bigr), \\
%     \beta(x)
%     &=
%     \operatorname{Mean}_{h\in\mathcal{H}_{\mathrm{sel}}}
%     m_h(x).
% \end{aligned}
% \end{equation*}
% \begin{equation*}
%     \theta_h^{(-x)}
%     =
%     \frac{1}{|\mathcal{D}_{\mathrm{public}}|-1}
%     \sum_{y\in\mathcal{D}_{\mathrm{public}}\setminus\{x\}}m_h(y),
% \end{equation*}
% and define task discrimination and difficulty as

% Positive $\alpha(x)$ indicates that stronger harnesses perform better on $x$, whereas non-positive values identify tasks that provide little reliable discrimination. This measure is analogous to corrected item--total correlation and captures structural responsiveness more directly than raw score variance.

\paratitle{Task Suite Construction.}
We first remove tasks with non-positive $\mathrm{Sens}(x)$, as they provide limited signal for distinguishing harness quality. Among the remaining tasks, we partition them into difficulty strata according to $1-\mathrm{Perf}(x)$ and select tasks with the highest $\mathrm{Sens}(x)$ within each stratum. This procedure yields a task set that preserves difficulty diversity while prioritizing tasks that are more responsive to harness evolution. Then we randomly split tasks within each stratum into validation and evaluation suites, ensuring that they follow the same difficulty distribution. Table~\ref{tab:task-selection-stats} reports the intermediate statistics of the task selection process.


% Within each source, we discard negatively discriminative tasks when quotas permit, stratify the remainder by difficulty, and favor harder strata to preserve performance headroom. We select the highest-discrimination tasks within each stratum and alternate them between the train and evaluation suites by discrimination rank, balancing responsiveness across the two splits. Table~\ref{tab:alpha-selection} reports the resulting statistics.
